\section{DFT 是变换，FFT 是算法}
\label{sec:08-Numerical-Analysis/08-Fourier-and-Spectral-Methods/03}

一秒时长的 CD 音质音频包含 \(N=48000\) 个采样点。想知道这段音频中 440Hz 的 A 音有多强，需要把它从时间序列变换到频率坐标。离散 Fourier 变换（DFT）的公式谁都会写——但直接算一遍的全部 \(N\) 个频率分量，求和项的数量级约 \(N^2\approx 2.3\times10^9\)。就是两亿三千个操作。

而快速 Fourier 变换（FFT）靠一个巧妙的递归分解，把同样的 DFT 结果用约 \(N\log_2 N\approx 7.5\times10^5\) 量级的操作算出来。减少的不是百分之几十，是三个数量级。这三个数量级的差异，直接决定了频谱显示能不能实时、卷积滤波能不能不卡顿。

本节把 DFT 这个数学变换和 FFT 这个计算算法拆开来，先讲 DFT 是什么（它不是一个近似，而是一个精确的基变换），再讲 FFT 是怎么用聪明办法算它的。

\subsection{有限维向量也有频率基}

长度为 \(N\) 的序列 \(x_0,\ldots,x_{N-1}\) 属于 \(\C^N\)——这就是一个 \(N\) 维复向量。离散 Fourier 变换把标准基 \(\{e_0,\ldots,e_{N-1}\}\) 旋转到频率基：
\[
X_k=\sum_{n=0}^{N-1} x_n\,e^{-2\pi i k n/N},\qquad k=0,\ldots,N-1.
\]

逆变换几乎一样，只差一个共轭和归一化：
\[
x_n=\frac1N\sum_{k=0}^{N-1} X_k\,e^{2\pi i k n/N}.
\]

这两组公式在说什么？每一个频率索引 \(k\) 对应一个复正弦波 \(e^{2\pi i k n/N}\)，它在 \(N\) 个采样点上以频率 \(k/N\)（单位：周每样本）振荡。\(X_k\) 就是这个频率分量的复数振幅。

为什么这组基是正交的？单位根做几何级数求和：
\[
\sum_{n=0}^{N-1} e^{2\pi i (k-\ell)n/N}
= \begin{cases}N, & k\equiv\ell\pmod N,\\0, & \text{否则}.\end{cases}
\]
当 \(k\neq\ell\) 时，这些复数均匀分布在单位圆上，向量和恰好为零——对称性使它们完全抵消。这个正交性是 DFT 能从 \(X_k\) 无损回到 \(x_n\) 的数学保证。

用 \(N=4\) 确认手感。四个单位根是 \(\omega^0=1,\omega^1=i,\omega^2=-1,\omega^3=-i\)（其中 \(\omega=e^{-2\pi i/4}=-i\)）。四个频率基向量为：
\[
\begin{aligned}
k=0&:\;(1,1,1,1) & \text{直流——全一}\\
k=1&:\;(1,\omega,\omega^2,\omega^3)=(1,-i,-1,i)\\
k=2&:\;(1,-1,1,-1) & \text{Nyquist——最高频}\\
k=3&:\;(1,i,-1,-i) & \text{负频——\(k=1\) 的共轭}
\end{aligned}
\]
任意两个不同 \(k\) 的频率基向量做内积，四个复数均匀分布，和为零。正交性让你可以分别读出每个频率分量的强度，而不受其他频率干扰。

\textbf{归一化约定不统一。}不同库把 \(1/N\) 放在正变换、逆变换，或两边各放 \(1/\sqrt{N}\)。比较不同来源的数值结果前，先确认归一化——否则差一个全局因子。

\subsection{DFT 隐含了周期延拓}

DFT 把长度为 \(N\) 的序列看作某个周期序列的完整一周期。索引 \(n\pm N\) 处的值被暗中与 \(n\) 处的值等同。这个假设不是后加的——它来自复指数 \(e^{-2\pi i k n/N}\) 对 \(n\) 的 \(N\)-周期性。由此产生三个直接后果。

\textbf{频率索引超过 \(N/2\) 是负频率。}\(X_{N-k}=X_{-k}\)。对实序列有共轭对称：
\[
X_{N-k}=\overline{X_k},
\]
所以只看 \(k=0,\ldots,\lfloor N/2\rfloor\) 就捕获了全部谱信息（直流、正频率、Nyquist 频率）。再往后是它们的镜像。

\textbf{频率分辨率由时长决定。}若采样间隔为 \(\Delta t\)，相邻频率 bin 的间距为
\[
\Delta\nu=\frac{1}{N\Delta t}.
\]
你只能分辨频率差至少为 \(1/(N\Delta t)\) 的两个正弦波。想看到更细的频谱结构？增加观测时长，而不是提高采样率。

\textbf{DFT 只描述窗口内信号的周期坐标。}它不会告诉你采样前信号有没有混叠——那一步失误发生在连续域，不在 DFT 内。它也不会帮你区分信号的自然周期和你人为截断产生的泄漏——那是窗函数的选择问题。

\subsection{圆周卷积从周期性自然出现}

DFT 下的逐点相乘对应时域圆周卷积：
\[
(x \circledast y)_n = \sum_{m=0}^{N-1} x_m y_{(n-m)\bmod N}.
\]
索引越过 \(N-1\) 后绕回开头——因为 DFT 把序列看成循环的。

如果你想用 FFT 高效计算两个序列的\textbf{线性}卷积（生活中的正常卷积），圆周卷绕会污染结果：尾部绕回头部，产生 circular aliasing。解法是把两个序列零填充到至少 \(N+M-1\) 的长度再变换，让圆周卷绕发生在填充的零区域，不影响实际数据。

零填充在这里的用途是防 aliasing，和另一种常见用途——频谱绘图时的插值致密——完全不同。后者是在现有数据窗内加密采样 DTFT，不会增加频率分辨率。把两种零填充混为一谈是常见的踩坑点。

看一个具体的小例子来感受 circular aliasing。序列 \(a=[1,2,3]\)（长度 3）与 \(b=[4,5]\)（长度 2）的线性卷积应有 \(3+2-1=4\) 项。直接 FFT 长度 3 的圆周卷积会得到长度为 3 的结果，其中第三项被前两项的``尾巴''污染了。正确的做法是把二者都补零到长度 \(4\)（因为是 2 的幂，实际会补到 \(4\)），做完 FFT、逐点相乘、IFFT 后得到完整四元素结果。少补一个零，末尾的值就会绕回到开头。

\subsection{直接 DFT 为什么有大量重复工作}

固定 \(N\)，令 \(\omega_N=e^{-2\pi i/N}\) 为单位根。DFT 的 \(k\) 分量：
\[
X_k=\sum_{n=0}^{N-1} x_n \omega_N^{kn}.
\]
对每个 \(k\)，这需要 \(N\) 次复数乘法和 \(N-1\) 次加法。\(N\) 个分量总共约 \(O(N^2)\) 次运算。对 \(N=48000\)，就是约 23 亿次。

但 \(\omega_N^{kn}\) 的计算有大量冗余。\(\omega_N\) 是一个 \(N\) 次单位根，\(\omega_N^{kn}=\omega_N^{kn\bmod N}\)——只有 \(N\) 个不同的幂次。而且很多乘积在不同的 \(k\) 之间可以共享。

Cooley 和 Tukey 在 1965 年找到了一种重组计算的方式。当 \(N\) 为偶数时，把输入按偶索引和奇索引分组。利用 \(\omega_N^2=\omega_{N/2}\)：
\[
\begin{aligned}
X_k &= \sum_{m=0}^{N/2-1} x_{2m}\,\omega_{N/2}^{mk}
      + \omega_N^k \sum_{m=0}^{N/2-1} x_{2m+1}\,\omega_{N/2}^{mk}\\[4pt]
    &= E_k + \omega_N^k\,O_k.
\end{aligned}
\]

两个求和正是偶数子序列和奇数子序列各自的 \(N/2\) 点 DFT——记为 \(E_k\) 和 \(O_k\)。一个 \(N\) 点 DFT 被拆成两个 \(N/2\) 点 DFT 外加 \(O(N)\) 次组合。但还没完：因为 \(E_k\) 和 \(O_k\) 以 \(N/2\) 为周期，且 \(\omega_N^{k+N/2}=-\omega_N^k\)，后一半频点的值可以免费获得：
\[
X_{k+N/2}=E_k-\omega_N^k O_k.
\]

一个大小为 \(N\) 的问题分成两个大小为 \(N/2\) 的子问题，每个子问题又按同一方式继续拆分。递归深度 \(\log_2 N\) 层，每层合并代价 \(O(N)\)：
\[
T(N)=2T(N/2)+O(N)=O(N\log N).
\]

这就是 radix-2 Cooley--Tukey FFT。它不是近似，不是捷径——在精确算术下，FFT 的输出与直接 \(\sum_{n=0}^{N-1}x_n\omega_N^{kn}\) 完全一致。浮点下因运算次序不同会有微小的舍入差异，但它们算的是同一个数学量。

用 \(N=4\) 走通全过程。输入 \(x_0,x_1,x_2,x_3\)。
\begin{itemize}
\tightlist
\item 偶索引 DFT：\(E_0=x_0+x_2,\;E_1=x_0-x_2\)（大小 2 DFT）。
\item 奇索引 DFT：\(O_0=x_1+x_3,\;O_1=x_1-x_3\)。
\item 组合：
  \(X_0=E_0+\omega_4^0O_0=E_0+O_0\)，
  \(X_1=E_1+\omega_4^1O_1=E_1-iO_1\)，
  \(X_2=E_0+\omega_4^2O_0=E_0-O_0\)，
  \(X_3=E_1+\omega_4^3O_1=E_1+iO_1\)。
\end{itemize}
直接 DFT 需要 \(4\times4=16\) 次复数乘法（不利用对称性时），FFT 两层分解下来约 \(4\log_2 4=8\) 次——\(N=4\) 时优势不明显，但减半的逻辑已经在了。\(N\) 每翻倍，FFT 省下的倍数就更大。

\subsection{FFT 不是另一种变换}

FFT 是\textbf{算法族}——实现 DFT 的具体计算方案。长度不必是 2 的幂：混合 radix 处理复合数长度，Rader 和 Bluestein 算法覆盖素数长度。成熟的 FFT 库（FFTW 等）会在运行时根据长度自动选择最优分解计划。

把序列补零填到``快的长度''有时加快计算，但补零不会增添原始信息——频率分辨率的物理约束仍然由实际采样时长决定。

回到开头的数字。\(N=48000\) 时，\(N^2\approx2.3\times10^9\) 对 \(N\log_2N\approx7.5\times10^5\)，两者并非逐操作的精确比较（常数因子不同），但三个数量级的差距足以解释一件事：频谱分析、快速卷积、降噪滤波之所以能实时运行，靠的不是硬件堆砌，是算法从 \(O(N^2)\) 降到了近乎 \(O(N\log N)\)。

现实仍有权衡。一次处理更长的音频块提高频率分辨率和吞吐效率，但也增加等待整块数据到齐的延迟。复杂度让任务变成可能，但块长与窗函数的选择仍由应用的时间精度需求决定——FFT 给你速度，不给额外信息。

\subsection{使用 FFT 时要检查什么：一份排查清单}

傅里叶代码写完之后，第一个版本几乎必定有某个参数选错。以下是一份可以逐条核对的清单：

\begin{itemize}
\tightlist
\item
  变换轴和数据布局：数据是按行还是按列存放？多通道数据的变换轴是否正确？
\item
  实输入接口：库可能只返回非负频率的半谱（\(k=0,\ldots,\lfloor N/2\rfloor\)），长度不是 \(N\)。这个行为在 NumPy 的 \texttt{rfft} 和 \texttt{fft} 之间不同。
\item
  正逆归一化：正变换和逆变换之间是否满足 round-trip 恒等式 \(x = \operatorname{IFFT}(\operatorname{FFT}(x))\)？不同库把 \(1/N\) 放在不同位置。
\item
  频率单位和 bin 排列：是否需要 \texttt{fftshift} 把零频移到中心？频率轴是按 bin 索引还是 Hz 标注？
\item
  窗函数、去均值、去趋势：FFT 前是否去除了直流分量？是否应用了合适的窗函数来抑制频谱泄漏？
\item
  卷积时的零填充：长度是否至少为 \(N+M-1\)？补零不够会导致 circular aliasing 污染有效区域。
\item
  批量与复用：反复做同长度变换时，是否复用了计划（plan）和工作区以避免每次重新规划？
\end{itemize}

用几个已知信号快速验证回路。这四个测试值得背下来：
\begin{itemize}
\tightlist
\item 单位脉冲 \([1,0,0,\ldots,0]\)：DFT 后所有频率分量的模都应为 1（全一谱）——验证基向量的对称性。
\item 常数序列 \([c,c,c,\ldots,c]\)：DFT 后只有直流 \(k=0\) 处非零，值为 \(Nc\)——验证正交性的消除效果。
\item 恰好落在某个频率 bin 上的单频正弦：DFT 后应该仅在该 bin 和其共轭位置出现峰值——验证频率索引正确。
\item Forward--inverse 往返：IFFT(FFT(x)) 应等于原序列（在浮点精度内）——验证归一化因子和正逆变换的配对关系。
\end{itemize}
这四个测试能在几秒内揪出缩放和索引错误——比盯着代码找 bug 快得多。几乎所有 FFT 踩坑经验都可以归结为：先在一个你完全知道答案的信号上跑一遍，而不是在你真正关心的数据上猜。

\subsection{窗函数与频谱泄漏：DFT 不会帮你修的瑕疵}

DFT 把信号看作周期信号的一个周期段。但如果实际信号在这个窗口的起点和终点处不连续——比如你把一段正弦波在不是完整周期的地方截断了——DFT 看到的就不是一个干净的正弦波周期延拓，而是一个有跳变的拼接周期。这个跳变会向所有频率 bin 泄漏能量，让原本应该只有一个峰的频谱变成一片模糊的山脊。

窗函数（Hann、Hamming、Blackman 等）的作用是在做 DFT 之前，把信号两端逐渐衰减到零，消除端点不连续。代价是主瓣变宽——频率分辨率下降。不加窗的矩形窗分辨率最高但泄漏最严重；窗越平滑，泄漏越少但分辨率越粗。

选择窗函数时面对的是一个权衡：你想要准确的高度（低泄漏），还是准确的位置（窄主瓣）？这个权衡不由 DFT/FFT 决定，而由你的信号长度和你关心的频率特征决定。FFT 给你速度和精确的计算结果——但它不会替你判断这个结果的质量。质量的上限在你选择窗口、决定截断长度的那一刻就锁定了。

\subsection{DFT 与连续 Fourier 级数：采样的代价}

DFT 是对连续信号采样后得到的离散工具。连续 Fourier 级数把一个周期函数分解为无穷个正弦波；DFT 把这个无穷级数截断成 N 个项。截断和采样各引入一个效应。

采样导致混叠：连续域中频率高于 Nyquist 频率 \(f_s/2\)（\(f_s=1/\Delta t\) 为采样率）的分量，在 DFT 谱中折叠到低频位置，与真实低频分不开。混叠在采样前就发生了——DFT 只是忠实地报告了被混叠后的结果。

截断导致泄漏：有限长度的数据窗口等价于给信号乘了一个矩形窗。矩形窗的傅里叶变换是 sinc 函数，它在频域中有波动旁瓣。即使信号本身只有一个精确的频率，它的 DFT 谱也会在目标频率周围散开一片——这就是频谱泄漏。加窗函数是用一个更平滑的截断来换取更少的旁瓣泄漏。

这两个效应都不由 DFT 改正，也都不能靠 FFT 加速来弥补。它们是信号从连续到离散这条路上必须付出的代价。DFT 给你的是窗口内采样的精确频率坐标，不会多给你信息（混叠已经不可逆），也不会替你美化信息（泄漏需要窗函数处理）。理解 DFT 的边界，和使用 FFT 的速度，是同一枚硬币的两面。

\subsection{DFT 与 FFT：一个负责含义，一个负责速度}

DFT 是 N 维复向量空间中一次精确的基变换——把时间序列映射为 N 个复正弦波的叠加系数。它隐含周期延拓，伴随圆周卷积，频率分辨率由观测时长锁定。这些性质不是 DFT 的 bug，而是它的定义。

FFT 是把这个数学定义的输出以 \(O(N\log N)\) 算出来的工程方案。在精确算术下，FFT 和直接 DFT 算的是完全相同的 \(X_k\)。浮点下因运算次序不同，舍入误差的分布略有不同，但差异的量级远小于应用层关心的精度范围。

把它俩分开理解：变换告诉你能做什么（把时间序列转到频率坐标、在频率坐标中高效卷积、通过对谱的操作实现滤波），算法告诉你这件事能在什么规模上做（毫秒级处理数万个采样点）。DFT 是数学，FFT 是工程——数学定义了极限，工程让极限不再挡路。
